% sec2_11.tex — Section 2.11: Application: Rational Expectations Econometrics

\section{Application: Rational Expectations Econometrics}
\label{sec:2.11}

In dynamic economic models and other contexts, expectations about the future
naturally play an important role. \textbf{Rational expectations econometrics} concerns
estimating equations involving expectations. Expectations are usually not observable, but, if we assume that economic agents form expectations rationally, we can
overcome the problem of unobservable expectations using the techniques of this
chapter. The application we consider is Fama's \textbf{efficient market hypothesis}. In
his own words, ``An efficient capital market is a market that is efficient in processing
information\ldots. In an efficient market, prices `fully reflect' available information''
(Fama, 1976, p.~133). The words ``fully reflect'' can be formalized precisely. We
will do so for a particular capital market and test its implication.

\subsection*{The Efficient Market Hypotheses}

The capital market studied in Fama (1975) is the market for U.S.\ Treasury bills. In
this section we focus on the one-month Treasury bill rates observed on a monthly
basis. (Later in the book we will also consider Treasury bills of different maturities.) To formalize market efficiency for the bills market, we need to introduce a
few concepts from macroeconomics and finance. Since in this section we deal
exclusively with time-series data, we use ``$t$'' rather than ``$i$'' for the subscript.
Define
%
\begin{align*}
v_t &\equiv \text{price of a one-month Treasury bill at the beginning of month } t, \\
R_t &\equiv \text{one-month nominal interest rate over month } t, \text{ i.e., nominal return} \\
    &\quad \text{on the bill over the month} = (1 - v_t)/v_t,
      \text{ so } v_t = 1/(1 + R_t), \\
P_t &\equiv \text{value of the Consumer Price Index (CPI), which is our measure of} \\
    &\quad \text{the price level at the beginning of month } t, \\
\pi_{t+1} &\equiv \text{inflation rate over month } t
      \text{ (i.e., from the beginning of month } t \text{ to} \\
    &\quad \text{the beginning of month } t + 1)
      = (P_{t+1} - P_t)/P_t, \\
{}_t\pi_{t+1} &\equiv \text{expected inflation rate over month } t,
      \text{ expectation formed at the} \\
    &\quad \text{beginning of month } t, \\
\eta_{t+1} &\equiv \text{inflation forecast error} = \pi_{t+1} - {}_t\pi_{t+1}, \\
r_{t+1} &\equiv \text{ex-post real interest rate over month } t
      = \frac{1/P_{t+1} - v_t/P_t}{v_t/P_t} \\
    &= \frac{1 + R_t}{1 + \pi_{t+1}} - 1 \approx R_t - \pi_{t+1}, \\
{}_tr_{t+1} &\equiv \text{ex-ante real interest rate}
      = \frac{1 + R_t}{1 + {}_t\pi_{t+1}} - 1
      \approx R_t - {}_t\pi_{t+1}.
\end{align*}

% Figure 2.2 placeholder
\begin{figure}[htbp]
\centering
\includegraphics[width=0.9\textwidth,trim=50 350 50 350,clip]{figures/fig_2_2.png}
\caption{What's Observed When}
\label{fig:2.2}
\end{figure}

Timing---the value of which variable is observed at what point in time---is very
important in rational expectations econometrics. Our rule for dating variables
(which is different from Fama's) is that the variable has subscript $t$ if its value
is first observed at the beginning of period (month) $t$.\footnote{In Fama's (1975)
notation, our $\pi_{t+1}$ is $-\widetilde{\Delta}_t$, $r_{t+1}$ is $\tilde{r}_{t+1}$,
and $I_t$ is $\phi_t$.}
For example, because $\pi_{t+1}$,
the inflation rate over month $t$, depends on $P_{t+1}$, its subscript is $t+1$ rather than $t$.
For the same reason, $r_{t+1}$, the ex-post real interest rate over month $t$ has subscript
$t + 1$. In contrast, the nominal interest rate $R_t$ over month $t$ has subscript $t$, not
$t + 1$, because it is determined by the price of the Treasury bill at the beginning of
month $t$, $v_t$. Figure~\ref{fig:2.2} shows the value of which variable is observed at what point
in time.

We take a period to be a month, only because our data are monthly. Note that
the maturity of the security (one month) coincides with the sampling interval (a
month). If the sampling interval is finer than the maturity, which is the case if,
for example, we have monthly data on the three-month Treasury bill rate, then
maturities overlap and we need a more sophisticated technique to be introduced in
Chapter~6.

The efficient market hypothesis is a joint hypothesis combining:

\medskip
\noindent\textbf{Rational Expectations.}\quad
Inflationary expectations are rational:
${}_t\pi_{t+1} = \E(\pi_{t+1} \mid I_t)$, where $I_t$ is information available at the beginning of month
$t$ and includes $\{R_t, R_{t-1}, \ldots, \pi_t, \pi_{t-1}, \ldots\}$.
Also, $I_t \supseteq I_{t-1} \supseteq I_{t-2} \supseteq \cdots$.
That is, agents participating in the market do not forget.

\medskip
\noindent\textbf{Constant Real Rates.}\quad
The ex-ante real interest rate is constant: ${}_tr_{t+1} = r$.

\subsection*{Testable Implications}

We derive two testable implications of the efficient market hypothesis by utilizing the following key observations about the inflation forecast error under rational
expectations.

\begin{enumerate}[label=(\alph*)]
\item $\E(\eta_{t+1} \mid I_t) = 0$, that is, the inflation forecast error is a martingale difference
sequence with respect to the information set.

\item $\{\eta_t, \eta_{t-1}, \ldots\}$ as well as
$\{R_t, R_{t-1}, \ldots, \pi_t, \pi_{t-1}, \ldots\}$ are included in $I_t$ (known
at the beginning of month $t$). That is, agents remember all past mistakes.
\end{enumerate}

(a) makes intuitive sense; if people use all available information to forecast the
inflation rate, the forecast error, which is known only after the fact, will not have
any systematic relation to what people knew when they formed expectations. It can
be proved easily:
%
\begin{align}
\E(\eta_{t+1} \mid I_t)
&= \E(\pi_{t+1} - {}_t\pi_{t+1} \mid I_t)
  \quad (\text{since } \eta_{t+1} \equiv \pi_{t+1} - {}_t\pi_{t+1}) \notag \\
&= \E(\pi_{t+1} \mid I_t) - \E({}_t\pi_{t+1} \mid I_t) \notag \\
&= \E(\pi_{t+1} \mid I_t) - \E[\E(\pi_{t+1} \mid I_t) \mid I_t] \notag \\
&\quad (\text{since } {}_t\pi_{t+1} = \E(\pi_{t+1} \mid I_t)
  \text{ by rational expectations}) \notag \\
&= \E(\pi_{t+1} \mid I_t) - \E(\pi_{t+1} \mid I_t) = 0.
\label{eq:2.11.1}
\end{align}
%
(b) follows because agents remember past inflation forecasts. Since
${}_t{}_{-j-1}\pi_{t-j} = \E(\pi_{t-j} \mid I_{t-j-1})$ is a function of $I_{t-j-1}$, it is included in
$I_{t-j-1}$ (i.e., known at the beginning of month $t - j - 1$) and hence in $I_t$ (known at the beginning of
month $t$) for $j \geq 0$. Thus
$\eta_{t-j} \equiv \pi_{t-j} - {}_{t-j-1}\pi_{t-j}$ is included in $I_t$ for all $j \geq 0$.
Observations (a) and (b), together with the Law of Total Expectations, imply

\begin{enumerate}[label=(\alph*),start=3]
\item $\{\eta_t\}$ is m.d.s.
\end{enumerate}

So, it is a zero mean and serially uncorrelated process.

These are not testable because expectations are unobservable, but, when combined with the constant-real-rate assumption, they imply two testable implications
about the inflation rate and the interest rate which are observable.

\medskip
\noindent\textbf{Implication 1:}\quad
The ex-post real interest rate has a constant mean and is serially
uncorrelated.

\medskip
This follows because
%
\begin{align}
r_{t+1} &\equiv R_t - \pi_{t+1} \notag \\
&= (R_t - {}_t\pi_{t+1}) - (\pi_{t+1} - {}_t\pi_{t+1}) \notag \\
&= {}_tr_{t+1} - \eta_{t+1} \quad \text{(by definition)} \notag \\
&= r - \eta_{t+1} \quad
  (\text{since } {}_tr_{t+1} = r \text{ by assumption}).
\label{eq:2.11.2}
\end{align}
%
By (c), $\{r_t\}$ has mean $r$ and is serially uncorrelated.

\medskip
\noindent\textbf{Implication 2:}\quad
$\E(\pi_{t+1} \mid I_t) = -r + R_t$.

\medskip
That is, the best (in the sense of minimum mean squared error) inflation forecast on
the basis of all the available information is the nominal interest rate; the nominal
interest rate $R_t$, which is determined by the price of a Treasury bill $v_t$, summarizes all the
currently available information relevant for predicting future inflation. This is the
formalization of asset prices ``fully reflecting'' available information. This, too, can
be derived easily. Solve \eqref{eq:2.11.2} for $\pi_{t+1}$ as
$\pi_{t+1} = -r + R_t + \eta_{t+1}$. So
%
\begin{align}
\E(\pi_{t+1} \mid I_t)
&= \E(-r + R_t + \eta_{t+1} \mid I_t) \notag \\
&= -r + R_t + \E(\eta_{t+1} \mid I_t)
  \quad (\text{since } r \text{ is a constant and } R_t \text{ is included in } I_t) \notag \\
&= -r + R_t \quad \text{(by (a))}.
\label{eq:2.11.3}
\end{align}
%
In the rest of this section, we test these two implications of market efficiency.

\subsection*{Testing for Serial Correlation}

Consider first the implication that the ex-post real rate has no serial correlation,
which Fama tests using the result from Proposition~\ref{prop:2.9}. We use the monthly data
set used by Mishkin (1992) on the one-month $T$-bill rate and the monthly CPI inflation rate, stated in percent at annual rates.\footnote{The data set is described in some
detail in the empirical exercise. Following Fama, the inflation measure used in the
text is calculated from the raw CPI numbers. The inflation measure used by Mishkin (1992)
treats the housing component of the CPI consistently. See question~(k) of Empirical
Exercise~1.}
Figure~2.3 plots the data. The ex-post
real interest rate, defined as the difference between the T-bill rate and the CPI inflation rate, is plotted in Figure~2.4. To duplicate Fama's results, we take the sample
period to be the same as in Fama (1975), which is January 1953 through July 1971.
The sample size is thus 223. The time-series properties of the real interest rate are
summarized in Table~\ref{tab:2.1}. In calculating the sample autocorrelations
$\hat{\rho}_j = \hat{\gamma}_j/\hat{\gamma}_0$,

% Figure 2.3 placeholder
\begin{figure}[htbp]
\centering
\includegraphics[width=0.9\textwidth,trim=50 350 50 350,clip]{figures/fig_2_3.png}
\caption{Inflation and Interest Rates}
\label{fig:2.3}
\end{figure}

% Figure 2.4 placeholder
\begin{figure}[htbp]
\centering
\includegraphics[width=0.9\textwidth,trim=50 350 50 350,clip]{figures/fig_2_4.png}
\caption{Real Interest Rates}
\label{fig:2.4}
\end{figure}

% Table 2.1
\begin{table}[htbp]
\centering
\caption{Real Interest Rates, January 1953--July 1971}
\label{tab:2.1}
\smallskip
\begin{tabular}{l*{12}{r}}
\toprule
\multicolumn{13}{c}{mean $= 0.82\%$, standard deviation $= 2.847\%$,
sample size $= 223$} \\
\midrule
& $j=1$ & $j=2$ & $j=3$ & $j=4$ & $j=5$ & $j=6$
& $j=7$ & $j=8$ & $j=9$ & $j=10$ & $j=11$ & $j=12$ \\
\midrule
$\hat{\rho}_j$
& $-0.101$ & $0.172$ & $-0.019$ & $-0.004$ & $-0.064$
& $-0.021$ & $-0.092$ & $0.095$ & $0.094$ & $0.019$
& $0.004$ & $0.207$ \\
std.\ error
& $0.067$ & $0.067$ & $0.067$ & $0.067$ & $0.067$
& $0.067$ & $0.067$ & $0.067$ & $0.067$ & $0.067$
& $0.067$ & $0.067$ \\
Ljung--Box $Q$
& $2.3$ & $9.1$ & $9.1$ & $9.1$ & $10.1$
& $10.2$ & $12.1$ & $14.2$ & $16.3$ & $16.4$
& $16.4$ & $26.5$ \\
$p$-value (\%)
& $12.8\%$ & $1.1\%$ & $2.8\%$ & $5.8\%$ & $7.3\%$
& $11.7\%$ & $9.6\%$ & $7.6\%$ & $6.1\%$ & $8.9\%$
& $12.8\%$ & $0.9\%$ \\
\bottomrule
\end{tabular}
\end{table}

\noindent
we use the formula \eqref{eq:2.10.1} for $\hat{\gamma}_j$ where the denominator is the sample size rather
than $n - j$. Similarly, the formula for the standard error of $\hat{\rho}_j$ is $1/\sqrt{n}$ rather
than $1/\sqrt{n-j}$ (so the standard error does not depend on $j$). The Ljung--Box $Q$ statistic
in the table for, say, $j = 2$ is \eqref{eq:2.10.5} with $p = 2$. The results in the table show
that none of the autocorrelations are highly significant individually or as a group,
in accordance with the first implication of the efficient market hypothesis.

\subsection*{Is the Nominal Interest Rate the Optimal Predictor?}

We can test Implication~2 by estimating an appropriate regression and carrying
out the $t$-test of Proposition~2.3. The remarkable fact about the efficient market
hypothesis is that, despite its simplicity, it is specific enough to imply the important
parts of the assumptions justifying the $t$-test.\footnote{The technical
part---that the fourth-moment matrix $\E(\mathbf{g}_t\mathbf{g}_t')$ exists and is finite and that the
regressors have finite fourth moments---is not implied by the efficient market hypothesis
and needs to be assumed.}
To see this, let $y_t \equiv \pi_{t+1}$,
$\mathbf{x}_t \equiv (1, R_t)'$, $\varepsilon_t \equiv \eta_{t+1}$ (the inflation forecast error), and rewrite
\eqref{eq:2.11.2} as
%
\begin{equation}
\pi_{t+1} = \beta_1 + \beta_2 R_t + \eta_{t+1}
\quad \text{or} \quad
y_t = \mathbf{x}_t'\bm{\beta} + \varepsilon_t,
\label{eq:2.11.4}
\end{equation}
%
so Assumption~2.1 is obviously satisfied with
%
\begin{equation}
\bm{\beta} = (-r, 1)'.
\label{eq:2.11.5}
\end{equation}
%
The other assumptions of the model are verified as follows.

\begin{itemize}
\item \emph{Assumption~2.3 (predetermined regressors).}\quad
The $\mathbf{g}_t$ ($\equiv \mathbf{x}_t \cdot \varepsilon_t$) in the present
context is $(\eta_{t+1}, R_t\eta_{t+1})'$. What we need to show is that
$\E(\eta_{t+1}) = 0$ and $\E(R_t\eta_{t+1}) = 0$. The former is implied by (a) and the Law of Total Expectations:
%
\begin{equation}
\E(\eta_{t+1}) = \E[\E(\eta_{t+1} \mid I_t)] = 0.
\label{eq:2.11.6}
\end{equation}
%
The latter holds because
%
\begin{align}
\E(R_t\eta_{t+1})
&= \E[\E(R_t\eta_{t+1} \mid I_t)]
  \quad \text{(by the Law of Total Expectations)} \notag \\
&= \E[R_t\,\E(\eta_{t+1} \mid I_t)]
  \quad \text{(by linearity of conditional expectations and }
  R_t \in I_t\text{)} \notag \\
&= 0 \quad \text{(by (a))}.
\label{eq:2.11.7}
\end{align}

\item \emph{Assumption~2.5.}\quad
By (b), $I_t$ includes $(\varepsilon_{t-1}\ (= \eta_t), \varepsilon_{t-2}, \ldots)$ as well as
$(\mathbf{x}_t, \mathbf{x}_{t-1}, \ldots)$. (a) then implies (2.3.1), which is a sufficient condition for
$\{\mathbf{g}_t\}$ to be m.d.s.

\item \emph{Assumption~2.4.}\quad
$\E(\mathbf{x}_t\mathbf{x}_t')$ here is
%
\begin{equation}
\E(\mathbf{x}_t\mathbf{x}_t') = \begin{bmatrix}
1 & \E(R_t) \\ \E(R_t) & \E(R_t^2)
\end{bmatrix}.
\label{eq:2.11.8}
\end{equation}
%
Its determinant is $\E(R_t^2) - [\E(R_t)]^2 = \Var(R_t)$, which must be positive; if the
variance were zero, we would not observe the interest rate $R_t$ fluctuating.

\item \emph{Assumption~2.2 (ergodic stationarity).}\quad
It requires $\{\pi_{t+1}, R_t\}$ to be ergodic stationary. The plot in Figure~2.3 shows a stretch of upward movements, especially
for the nominal interest rate, although the stretch is followed by a reversion to a
lower level. Testing stationarity of the series will be covered in Example~9.2 of
Chapter~9. For now, despite this rather casual evidence to the contrary, we proceed under the assumption of stationarity (but the implication of having trending
series will be touched upon in the final subsection). For ergodicity, we just have
to assume it.
\end{itemize}

What we have shown is that, under the efficient market hypothesis, $\{y_t, \mathbf{x}_t\}$
(where $y_t = \pi_{t+1}$ and $\mathbf{x}_t = (1, R_t)'$) belongs to the set of DGPs satisfying
Assumptions~2.1--2.5.

Having verified the required assumptions, we now proceed to estimation. Using
our data, the estimated regression is, for $t = 1/53, \ldots, 7/71$,
%
\begin{equation}
\begin{aligned}
\pi_{t+1} &= \underset{(0.431)}{-0.868} + \underset{(0.112)}{1.015}\; R_t, \\[4pt]
R^2 &= 0.24, \quad
\text{mean of dependent variable} = 2.35, \quad
\mathit{SER} = 2.84\%, \quad
n = 223.
\end{aligned}
\label{eq:2.11.9}
\end{equation}
%
The heteroskedasticity-robust standard errors calculated as in (2.4.1) (where
$\widehat{\mathbf{S}}$ is given by (2.5.1) with $\hat{\varepsilon}_i = e_i$) are shown in parentheses.
As shown in \eqref{eq:2.11.5}, market efficiency implies that the $R_t$ coefficient equals
1.\footnote{The intercept is known to be $-r$, but it does not produce a testable restriction
since the efficient market hypothesis does not specify the value of $r$.}
The robust $t$-ratio for the null that the $R_t$ coefficient equals 1 is
$(1.015 - 1)/0.112 = 0.13$. Its $p$-value is about 90 percent. Thus, we can easily
accept the null hypothesis.

There are other ways to test the implication of the nominal interest rate summarizing all the currently available information. We can bring in more explanatory
variables in the regression and see if they have an explanatory power over and
above the nominal rate. The additional variables to be brought in, however, must
be part of the current information set $I_t$ because, obviously, the efficient market

\subsection*{$R_t$ Is Not Strictly Exogenous}

As has been emphasized, one important advantage of large-sample theory over
finite-sample theory of Chapter~1 is that the regressors do not have to be strictly
exogenous as long as they are orthogonal to the error term. We now show by
example that $R_t$ is not strictly exogenous, so finite-sample theory is not applicable
to the Fama regression \eqref{eq:2.11.4}. Suppose that the process for the inflation rate is
an AR(1) process:
\[
\pi_t = c + \rho\pi_{t-1} + \eta_t, \quad
\{\eta_t\} \text{ is independent white noise.}
\]
If $\eta_{t+1}$ is independent of any elements of $I_t$, then
\begin{align*}
\E(\pi_{t+1} \mid I_t)
&= c + \rho\pi_t + \E(\eta_{t+1} \mid I_t)
  \quad (\text{since } \pi_t \text{ is in } I_t) \\
&= c + \rho\pi_t
  \quad (\text{since } \E(\eta_{t+1} \mid I_t) = \E(\eta_{t+1}) = 0 \text{ by assumption}).
\end{align*}
%
So $\eta_{t+1}$ is indeed the inflation forecast error for $\pi_{t+1}$ and
\[
R_t = r + c + \rho\pi_t
\]
(since $R_t = r + \E(\pi_{t+1} \mid I_t)$ under market efficiency). It is then easy to see that
the error term $\varepsilon_t$ in the Fama regression (of $\pi_{t+1}$ on a constant and $R_t$), which is the
inflation forecast error $\eta_{t+1}$, can be correlated with future regressors. For example,
%
\begin{align*}
\Cov(\varepsilon_t, R_{t+1})
&= \Cov(\eta_{t+1}, r + c + \rho\pi_{t+1}) \\
&= \rho\,\Cov(\eta_{t+1}, \pi_{t+1}) \\
&= \rho\,\Var(\eta_{t+1})
  \quad (\text{since } \pi_{t+1} = c + \rho\pi_t + \eta_{t+1}
  \text{ and } \Cov(\eta_{t+1}, \pi_t) = 0),
\end{align*}
%
which is not zero.

Although finite-sample theory is not applicable, its prescription for statistical
inference is approximately valid, provided that the error is conditionally homoskedastic and the sample size is large enough, which was the message of Section~2.6.

\subsection*{Subsequent Developments}

Although Fama's paper represents perhaps the finest example of rational expectations econometrics, events after its publication and a large number of empirical
studies it inspired proved that the near one-for-one relation between inflation and
interest rates is limited to the postwar period until 1979. Such strong association
cannot be found for the prewar period or for many other countries. Furthermore, the
increase in the real interest rates after the change in the Fed's operating procedure
that took place in October 1979 runs counter to the premise of the Fama regression that the expected real interest rate is constant over time. When the sample is
restricted to the post-October 1979 period, the interest rate coefficient in the Fama
regression is far below unity.\footnote{See, for example, Mishkin (1992, Table~1).
We will verify this in part~(l) of Empirical Exercise~1.}

These findings raise the question of why the strong association occurs only for
certain periods but not for others. The explanation by Mishkin (1992) is that the
inflation premium gets incorporated into interest rates only gradually. In periods
when the inflation rate shows only short-run fluctuations, interest rates are not
responsive to inflation. However, sustained movements in inflation will be reflected
in interest rates. The period when the strong association is observed (which
includes Fama's sample period, see Figure~2.3) is precisely when inflation showed
a sustained upward movement.

On January 29, 1997, the U.S.\ Treasury auctioned \$7 billion in ten-year
inflation-indexed bonds, making it possible for researchers to observe the ex-ante
real interest rate as the yield on indexed bonds. Evidence from Great Britain,
where indexed bonds have been available since the early 1980s, is that the yields
on indexed bonds, although much less volatile than the yield on ordinary bonds,
are not constant over time. The constant-real-rate assumption that we made was
the auxiliary assumption whereby we can test whether bond prices fully reflect
available information. Now we can do so without the auxiliary assumption by
regressing the actual inflation rate on the yield differential between the conventional and indexed bonds. We can also relax the assumption, made implicitly so
far, that the yield differential equals the expected inflation rate. Investors may not
be risk-neutral, and the inflation risk premium over and above the expected inflation rate may be needed to induce them to hold ordinary bonds. If so, the expected
inflation rate is once again unobservable. Rational expectations econometrics will
continue to be a useful tool for dealing with unobservable expectations.

\bigskip
\noindent\textsc{Questions for Review}
\begin{enumerate}
\item Let $R_t = 5\%$ and $\pi_{t+1} = 2\%$. Calculate $r_{t+1}$ using the two formulas,
$(1 + R_t)/(1 + \pi_{t+1}) - 1$ and $R_t - \pi_{t+1}$. Is the difference small? What if
$R_t = 20\%$ and $\pi_{t+1} = 17\%$?

\item Show that $\widehat{\E}^*(\pi_{t+1} \mid 1, R_t) = -r + R_t$ under the efficient market
hypothesis. \textbf{Hint:} Use \eqref{eq:2.9.7}.

\item If the inflation forecast error $\eta_{t+1}$ were observable, how would you test market
efficiency? Do we need the constant ex-ante real rate assumption?

\item If the inflation rates and the interest rates were measured in fractions rather
than percent (e.g., 0.08 instead of 8\%), how would the regression result \eqref{eq:2.11.9}
change? If the inflation rate is in percent but per month and the interest rate is
in percent per year, then $1 + y = (1 + x)^{12}$ so that $y \approx 12x$.
\textbf{Hint:} If $x$ is the inflation rate per month and $y$ is its annual
rate, then $1 + y = (1 + x)^{12}$ so that $y \approx 12x$.

\item Suppose in \eqref{eq:2.11.4} $\mathbf{x}_t$ includes a third variable which is not in $I_t$. Which part
of the argument deriving Assumption~2.3 fails? \textbf{Hint:} The third element of
$\mathbf{g}_t$ is the third variable times $\eta_{t+1}$.

\item Provide an example of market \emph{inefficiency} such that Implication~1 holds but
Implication~2 does not. \textbf{Hint:} Suppose $\pi_t$ and $R_t$ are serially independent and
mutually independent processes.
\end{enumerate}
